\subsection{主理想整环上挠模的典范分解}

设 M 是交换环 A 上的一个模。对于每个 \(\alpha\in A\) ，令 \(\mathbf{M}(\alpha)\) 表示 M 的自同态 \(x\mapsto\alpha x\) 的核。如果 \(\alpha\) 与 \(\beta\) 是 A 的两个元素，使得 \(\alpha\) 整除 \(\beta\) ，那么显然 \(\mathbf{M}(\alpha)\subset\mathbf{M}(\beta)\) 。特别地，当 n 遍历 \(\geqslant1\) 的有理整数集时，子模 \(\mathbf{M}(\alpha^{n})\) 构成一个递增序列；并集 \(M_{\alpha}\) （诸 \(\mathbf{M}(\alpha^{n})\) 的）因此是 M 的一个子模，由 M 中被 \(\alpha\) 的某个幂零化的元素组成。对于 M 的每个子模 N，显然有 \(N_{\alpha}=N\cap M_{\alpha}\) .

定义 1. --- 设 \(\pi\) 是主理想整环 A 的一个不可约元素；一个 A-模 M 被称为 \(\pi\) -准素的，如果对所有 \(x \in \mathbf{M}\) ，存在一个整数 \(n \geqslant 1\) 使得 \(\pi^n x = 0\) （换言之，如果 M 等于子模 \(\mathbf{M}_{\pi}\) ).

显然，每个形如 \(\mathbf{A}/(\pi^{s})\) 的循环模都是 \(\pi\) -准素的。对于任意一个 A-模 M，子模 \(M_{\pi}\) 是 \(\pi\) -准素的。

引理 1. --- 设 M 是主理想整环 A 上的一个模；对所有 \(\alpha \in \mathbf{A}\) （满足 \(\alpha \neq 0\) ），设 \(\alpha = \varepsilon \sum_{i=1}^{r} \pi_i^{n(i)}\) 是 \(\alpha\) 分解为不可约因子的一个分解（VII，第 4 页）。子模 \(\mathbf{N} = \mathbf{M}(\alpha)\) （由 M 中被 \(\alpha\) 零化的元素组成）是诸子模 \(\mathbf{M}(\pi_i^{n(i)})\) 的直和，且把每个 \(x \in \mathbf{M}(\alpha)\) 映到它在 \(\mathbf{M}(\pi_i^{n(i)})\) 中的分量的映射具有形式 \(x \mapsto \gamma_i x\) ( \(\gamma_i \in \mathbf{A}\) ）。此外

\[
\mathbf {M} \left(\pi_ {i} ^ {n (i)}\right) = \mathbf {N} \cap \mathbf {M} _ {\pi_ {i}} = \mathbf {N} _ {\pi_ {i}}.
\]

首先注意，N 被 \(\alpha\) 零化，因此具有自然的 \(\mathbf{A}/(\alpha)\) -模结构。根据 VII 第 3 页命题 4，从 \(\mathbf{A}/(\alpha)\) 到环的乘积 \(\mathbf{A}/(\pi_{i}^{n(i)})\) 的典范同态是环同构；现在应用 VII 第 6 页命题 1，我们推出 N 是 \(\mathbf{M}(\pi_{i}^{n(i)})\) 的直和；该分解的投影算子是位似变换，依据 VII 第 7 页注记 2。包含关系 \(\mathbf{M}(\pi_{i}^{n(i)})\subset\mathbf{M}(\alpha)\cap\mathbf{M}_{\pi_{i}}\) 是显然的；反之，设 \(x\in\mathbf{M}(\alpha)\cap\mathbf{M}_{\pi_{i}}\) 于是存在一个幂 \(\pi_{i}^{s}\) （ \(\pi_{i}\) 的幂）零化 x；我们可以假设 \(s\geqslant n(i)\) ；根据贝祖恒等式，存在 \(\lambda,\mu\in\mathbf{A}\) 使得 \(\pi_{i}^{n(i)}=\lambda\pi_{i}^{s}+\mu\alpha\) ，所以 \(\pi_{i}^{n(i)}x=0\) ，最终得 \(x\in\mathbf{M}(\pi_{i}^{n(i)})\) .

引理 2. —— 设 M 是整环 A 上的挠模（II，第 313 页）。对于 M 中元素的每个有限族 \((x_{i})_{1 \leq i \leq n}\) ，存在 A 中的一个元素 \(\gamma \neq 0\) 使得 \(x_{i}\) 都属于 \(\mathrm{M}(\gamma)\) .

事实上，对于每个指标 \(i\) ，在 A 中存在一个元素 \(\alpha_{i} \neq 0\) 零化 \(x_{i}\) ，且元素 \(\gamma = \prod_{i=1}^{n} \alpha_{i}\) 即符合要求。

定理1. --- 设 M 是主理想整环 A 上的挠模；对 A 的每个不可约元素 \(\pi\) ，令 \(M_{\pi}\) 为 M 的由被 \(\pi\) 的某个幂零化的元素组成的子模。若 P 是 A 的不可约元素的一个代表系，则 M 是其子模 \(M_{\pi}\) （ \(\pi \in P\) ）的直和.

每个元素 \(x \in M\) 都属于某个子模 \(\mathbf{M}(\alpha)\) （ \(\alpha \neq 0\) ），因此根据引理 1，它是有限个元素的和，其中每个元素都属于某个子模 \(M_{\pi}\) 。另一方面，如果 \(\sum_{\pi \in P} x_{\pi} = \sum_{\pi \in P} y_{\pi}\) ，其中 \(x_{\pi}, y_{\pi} \in M_{\pi}\) 对所有 \(\pi \in P\) 成立，且其中 \(x_{\pi}\) 与 \(y_{\pi}\) 除有限多个外均为零，则引理 2 表明存在 A 中的 \(\gamma \neq 0\) ，使得所有 \(x_{\pi}\) 与 \(y_{\pi}\) 都属于同一个子模 \(\mathbf{M}(\gamma)\) ；将引理 1 应用于 \(\mathbf{M}(\gamma)\) 表明 \(x_{\pi} = y_{\pi}\) 对所有 \(\pi \in P\) 成立，这就完成了证明。

显然，如果 \(\pi\) 与 \(\pi'\) 是两个相伴的不可约元素，则 \(M_{\pi} = M_{\pi'}\) ；因此，对于给定的模 M，子模 \(M_{\pi}\) 仅取决于 A 的理想 ( \(\pi\) )；它称为模 M 的 \(\pi\) -准素分量，而将 M 分解为诸 \(M_{\pi}\) 的直和称为 M 作为其 \(\pi\) -准素分量直和的典范分解。

推论1. --- 挠模 M 的每个子模 N 都是其子模 \(N \cap M_{\pi}\) 的直和.

这由以下事实得出： \(\mathbf{N} \cap \mathbf{M}_{\pi}\) 是 \(\pi\) -准素分量 \(\mathbf{N}_{\pi}\) （ \(\mathbf{N}\) 的）。

推论2.——挠 A-模 M 的子模 N 是直和因子，当且仅当 \(N_{\pi}\) 是 \(M_{\pi}\) （对 A 的每个不可约元素 \(\pi\) ）的直和因子。

确实，如果 \(\mathbf{N}\) 与 \(\mathbf{N}'\) 是 \(\mathbf{M}\) 的两个子模，则 \(\mathbf{M} = \mathbf{N} \oplus \mathbf{N}'\) 当且仅当 \(\mathbf{M}_{\pi} = \mathbf{N}_{\pi} \oplus \mathbf{N}_{\pi}'\) 对每个不可约元素 \(\pi\) （ \(\mathbf{A}\) 的）成立（推论 1）。

推论3.——设 N 是挠 A-模 M 的一个子模。如果对于 A 的每个不可约元素 \(\pi\) ，要么 \(N_{\pi} = 0\) ，要么 \((M/N)_{\pi} = 0\) ，则 N 是 M 的一个直因子。

事实上，条件 \((\mathbf{M}/\mathbf{N})_{\pi}=0\) 蕴含 \(N_{\pi}=M_{\pi}\) ，于是推论 2 适用。

一个 A-模 M 称为半单的，如果 M 的每个子模都是直因子（参见 A，VIII，§3）。

推论4. --- 设 A 是一个不是域的主理想整环，并设 M 是一个 A-模。则 M 是半单的，当且仅当 M 是挠模，且 \(\mathbf{M}_{\pi} = \mathbf{M}(\pi)\) 对 A 的每个不可约元素 \(\pi\) 成立。

首先假设 M 是半单的；令 \(x \in M\) 并且令 \(\pi\) 是 A 的一个不可约元素。若 N 是 \(A\pi x\) 在 M 中的一个补，那么我们可以写出 \(x = \alpha\pi x + y\) ，其中 \(\alpha \in A\) 且 \(y \in N\) ；但这意味着 \(y = (1 - \alpha\pi)x\) ，所以

\[
\pi (1 - \alpha \pi) x \in A \pi x \cap N = 0.
\]

首先由此得出 M 是一个挠模；如果此外 \(x \in M_{\pi}\) ，则 \(\pi(1 - \alpha\pi)x = 0\) ，从而 \(\pi x = \alpha\pi^{2}x = \alpha^{2}\pi^{3}x = \cdots = \alpha^{n}\pi^{n+1}x\) 为零，且 \(\mathbf{M}_{\pi} = \mathbf{M}(\pi)\) .

反之，由推论 2 可知，只需证明被一个不可约元素 \(\pi\) 零化的 A-模 M 是半单的；但这是显然的，因为此时 M 具有域 \(\mathrm{A}/(\pi)\) 上向量空间的自然结构，并且在此结构下，M 的子模恰好就是向量子空间。

注记 1. --- 显然，每个 \(\neq 0\) 元素（ \(\pi\) -准素模中的）的零化子都具有形式 \(A\pi^k (k > 0\) 为一个整数），因为它是包含 \(\pi\) 的一个幂的主理想。设 \(x\) 为 \(M\) 的一个元素；对每个 \(\pi \in P\) ，设 \(x_{\pi}\) 为 \(x\) 在 \(M_{\pi}\) 中的分量； \(x\) 的零化子是诸非零 \(x_{\pi}\) 的零化子的最小公倍元，但根据上述内容，在这种情形下它等于诸非零 \(x_{\pi}\) 的零化子的乘积（VI，第 16 页，命题 12（DIV））。

命题2. --- 若 M 是主理想整环 A 上的有限生成挠模，则 M 的 \(\pi\) -准素分量除有限个外均为零，且 M 在这些分量 \(M_{\pi}\) 上的投影算子是位似变换。

这由引理1直接得出，因为根据引理2，存在 \(\alpha \neq 0\) 在 A 中使得 \(\mathbf{M} = \mathbf{M}(\alpha)\) .

注2. --- 由 VII 第 7 页注记 3，一个有限生成的挠 A-模是循环的，当且仅当其每个 \(\pi\) -准素分量都是循环的。

定理1和命题2适用的一个重要特殊情形是当A=Z时；一个Z-模就是一个阿贝尔群。一个阿贝尔群被称为挠群，如果它是一个挠Z-模，也就是说它的所有元素都有有限阶。此时，取 P 为 \textgreater{} 0 的素数集合；对于每个素数 p \textgreater{} 0，一个（阿贝尔）群被称为 p-挠群，如果其所有元素的阶都是 p 的幂。在这个术语下，定理 1 表明每个挠阿贝尔群都是 p-挠群的直和。在有限群的情形下，这也由I，第80页，定理4得出。

